% Source/coverage: Weinberg QFT Vol. I, Ch. 12, Sec. 12.3, physical PDF
% pp. 539--548 (printed pp. 516--525), beginning at the section heading near
% the foot of p. 539 and ending immediately before the Sec. 12.4 heading.
% Inventory: equations (12.3.1)--(12.3.2); eight unnumbered displays;
% Table 12.2; Figures 12.4--12.5; three ordinary numbered footnotes.

\subsection{Is Renormalizability Necessary?}
\label{sec:12.3}

In the previous section we found a special class of theories having only
a finite number of terms in the Lagrangian, to which the renormalization
program is nevertheless applicable. These are theories in which all
interactions satisfy the renormalizability condition
\[
\Delta_i\equiv 4-d_i-\sum_f n_{if}(s_f+1)\geq 0,
\]
where $d_i$ and $n_{if}$ are the numbers of derivatives and fields of type
$f$ in interactions of type $i$, and $s_f$ is (with some qualifications) the
spin of fields of type $f$. For renormalization to work in such theories, it
is also usually necessary that all renormalizable interactions that are
allowed by symmetry principles should actually appear in the Lagrangian.

It is important that there are only a limited number of such interaction
types. $\Delta_i$ becomes negative if we have too many fields or derivatives,
or fields of too high spin. Barring special cancellations, there are no
renormalizable interactions at all involving fields with $s_f\geq 1$, because
the only possible term in the Lagrangian with $\Delta_i\geq 0$ that involves
such a field along with two or more other fields would involve a single
$s_f=1$ field along with two scalars and no derivatives, which would not be
Lorentz-invariant. We shall see in Volume II that general, massless, spin one
gauge fields in a suitable gauge effectively have $s_f=0$, like the photon.
Also, in Volume II we shall see that even massive gauge fields may effectively
have $s_f=0$, depending on where their mass comes from. Leaving aside these
special cases, Table~\ref{tab:weinberg-12-2} gives a list of \emph{all}
renormalizable terms in the Lagrangian density that are allowed by Lorentz
invariance and gauge invariance involving scalars ($s=0$), photons ($s=0$),
and spin $\tfrac12$ fermions ($s=\tfrac12$).

% Table 12.2, physical PDF p. 541 (printed p. 518).
\begingroup
\renewcommand{\thetable}{12.\arabic{table}}
\setcounter{table}{1}
\begin{table}[t]
\centering
\caption{Allowed renormalizable terms in a Lagrangian density involving
scalars $\phi$, Dirac fields $\psi$, and photon fields $A^\mu$. Here $n_{if}$
and $d_i$ are the number of fields of type $f$ and the number of derivatives
in an interaction of type $i$, and $\Delta_i$ is the dimensionality of the
associated coefficient.}
\label{tab:weinberg-12-2}
\renewcommand{\arraystretch}{1.15}
\setlength{\tabcolsep}{1.25em}
\begin{tabular}{ccc c c l}
\hline\hline
\multicolumn{3}{c}{$n_{if}$} & $d_i$ & $\Delta_i$ & $\mathcal H_i$ \\
Scalars & Photons & Spin $\tfrac12$ & & & \\
\hline
1 & 0 & 0 & 0 & 3 & $\phi$ \\
2 & 0 & 0 & 0 & 2 & $\phi^2$ \\
2 & 0 & 0 & 2 & 0 & $\partial_\mu\phi\partial^\mu\phi$ \\
3 & 0 & 0 & 0 & 1 & $\phi^3$ \\
4 & 0 & 0 & 0 & 0 & $\phi^4$ \\
2 & 1 & 0 & 1 & 0 & $\phi\partial_\mu\phi A^\mu$ \\
2 & 2 & 0 & 0 & 0 & $\phi^2 A_\mu A^\mu$ \\
1 & 0 & 2 & 0 & 0 & $\phi\bar\psi\psi$ \\
0 & 2 & 0 & 2 & 0 & $F_{\mu\nu}F^{\mu\nu}$ \\
0 & 0 & 2 & 0 & 1 & $\bar\psi\psi$ \\
0 & 0 & 2 & 1 & 0 & $\bar\psi\gamma^\mu\partial_\mu\psi$ \\
0 & 1 & 2 & 0 & 0 & $\bar\psi\gamma^\mu A_\mu\psi$ \\
\hline\hline
\end{tabular}
\end{table}
\endgroup


We see that the requirement of renormalizability puts severe restrictions on
the variety of physical theories that we may consider. Such restrictions
provide a valuable key to the structure of physical theories. For instance,
Lorentz and gauge invariance by themselves would allow the introduction of a
`Pauli' term proportional to $\bar\psi[\gamma_\mu,\gamma_\nu]\psi F^{\mu\nu}$
in the Lagrangian of quantum electrodynamics, which would make the magnetic
moment of the electron an adjustable parameter, but we exclude such terms
because they are not renormalizable. The successful predictions of quantum
electrodynamics, such as the calculation of the magnetic moment of the
electron outlined in Section 11.3, may be regarded as validations of the
principle of renormalizability. The same applies to the standard model of
weak, electromagnetic, and strong interactions, to be discussed in Volume
II; there are any number of terms that might be added to this theory, such as
four-fermion interactions among quarks and leptons, that would invalidate all
the predictions of the standard model, and are excluded only because they are
non-renormalizable.

Must we believe that the Lagrangian is restricted to contain only
renormalizable interactions? As we saw in the previous section, if we include
in the Lagrangian \emph{all} of the infinite number of interactions allowed by
symmetries, then there will be a counterterm available to cancel every
ultraviolet divergence. In this sense, as said earlier, non-renormalizable
theories are just as renormalizable as renormalizable theories, as long as we
include all possible terms in the Lagrangian.

In recent years it has become increasingly apparent that renormalizability is
not a fundamental physical requirement, and that in fact any realistic
quantum field theory will contain non-renormalizable as well as renormalizable
terms. This change in point of view can be traced in part to the continued
failure to find a renormalizable theory of gravitation. In the general class
of metric theories of gravitation governed by Einstein's principle of
equivalence there are no renormalizable interactions at all---generally
covariant interactions must be constructed from the curvature tensor and its
generally covariant derivatives, and hence, even in a `gauge' where the
graviton propagator goes as $k^{-2}$, these interactions involve too many
derivatives of the metric for renormalizability. In particular, we can easily
see that general relativity is non-renormalizable from the fact that its
coupling constant $8\pi G_N=(2.43\cdot 10^{18}\ \mathrm{GeV})^{-2}$ has
negative dimensionality. Even if nothing else did, the cancellation of
divergences due to virtual gravitons would require that the Lagrangian contain
all interactions allowed by symmetries---not only interactions involving
gravitons, but involving any particles.

But if renormalizability is not a fundamental physical principle, then how do
we explain the success of renormalizable theories like quantum
electrodynamics and the standard model? The answer can be seen by simple
dimensional analysis. We have already noted that the coupling constant of an
interaction of type $i$ has dimensionality
\begin{equation}
[g_i]\sim[\mathrm{mass}]^{\Delta_i},
\tag{12.3.1}\label{eq:12.3.1}
\end{equation}
where $\Delta_i$ is the index \eqref{eq:12.1.9}. Non-renormalizable
interactions are just those whose coupling constants have the dimensionality
of \emph{negative} powers of mass. Now, it is not unreasonable to guess from
\eqref{eq:12.3.1} that the coupling constants not only have dimensionalities
governed by $\Delta_i$, but are roughly of order
\begin{equation}
g_i\approx M^{\Delta_i},
\tag{12.3.2}\label{eq:12.3.2}
\end{equation}
where $M$ is some common mass. (This is found to be actually the case in the
effective field theories discussed below and in more detail in Volume II.) In
calculating physical processes at a characteristic momentum scale $k\ll M$,
the inclusion of a non-renormalizable interaction of type $i$ with
$\Delta_i<0$ will introduce a factor $g_i\approx M^{\Delta_i}$, which on
dimensional grounds must be accompanied by a factor $k^{-\Delta_i}$, and so
the effect of such an interaction is suppressed\footnote{It is essential at
this point to assume that the ultraviolet divergences have been removed by
renormalization, so that there are no factors of an ultraviolet cutoff
$\Lambda$ to mess up our dimensional analysis. Otherwise, dimensional
analysis tells us that for $\Lambda\to\infty$, each additional
non-renormalizable coupling constant factor $g_i$ with $\Delta_i<0$ would be
accompanied with a growing factor $\Lambda^{-\Delta_i}$. This dimensional
argument led Heisenberg~\hyperref[ref:12.9]{[9]} very early to classify
interactions according to the dimensionality of their coupling constants, and
to suggest~\hyperref[ref:12.10]{[10]} that new effects might arise at energies
of order $g_i^{1/\Delta_i}$, as for instance at the energy
$G_F^{-1/2}\approx 300\ \mathrm{GeV}$, where $G_F$ is the four-fermion coupling
constant of the Fermi beta decay theory. After the development of
renormalization theory it was noted by Sakata \emph{et
al.}~\hyperref[ref:12.11]{[11]} that the non-renormalizable theories are those
whose coupling constants have negative dimensionality.} for $k\ll M$ by a
factor $(k/M)^{-\Delta_i}\ll 1$. (This argument will be made more carefully
using the method of the renormalization group in Volume II.) The success of
the renormalizable theories of electroweak and strong energies shows only that
$M$ is very much larger than the energy scale at which these theories have
been tested.

For instance, the leading non-renormalizable corrections to the conventional
electrodynamics of electrons or muons would be those interactions of dimension
5, which are suppressed by only one factor of $1/M$. There is just one such
interaction allowed by Lorentz, gauge, and CP invariance, a Pauli term of order
$(ie/2M)\bar\psi[\gamma_\mu,\gamma_\nu]\psi F^{\mu\nu}$. According to
Eqs.~\eqref{eq:10.6.24}, \eqref{eq:10.6.17}, and
\eqref{eq:10.6.19}, such a term would contribute an
amount of order $4e/M$ to the magnetic moment of the electron or muon. The
calculated value of the magnetic moment of the electron agrees with experiment
to within terms of order $10^{-10}e/2m_e$, so $M$ must be greater than about
$8\cdot 10^{10}m_e=4\cdot 10^7\ \mathrm{GeV}$.

\begingroup
\emergencystretch=1em
This limit may be weakened if other symmetries restrict the form of the
non-renormalizable interactions. For instance, the conventional Lagrangian of
quantum electrodynamics is invariant under a chiral transformation
$\psi\to\gamma_5\psi$, except for a change of sign of the fermion mass term
$-m\bar\psi\psi$. If we assume that the full Lagrangian is invariant under a
formal symmetry $\psi\to\gamma_5\psi$, $m\to-m$, then a Pauli term in the
Lagrangian would have to appear with an extra factor $m/M$, so that its
contribution to the magnetic moment would be only of order $4em/M^2$. Because
of the extra factor of $m$, here it is the muon rather than the electron that
provides the most useful limit on $M$. The calculated value of the magnetic
moment of the muon agrees with experiment to within terms of order
$10^{-8}e/2m_e$, so $M$ must be greater than about
$\sqrt{8\cdot 10^8}\,m_\mu\allowbreak=3\cdot 10^3\ \mathrm{GeV}$. In any case, if $M$
is anywhere near as large as $10^{18}\ \mathrm{GeV}$, then we are certainly
justified in neglecting any non-renormalizable interactions that might appear
in quantum electrodynamics.\par
\endgroup

These considerations help us to cope with some of the puzzles associated with
higher-derivative terms in the Lagrangian. For instance, in the general theory
of a real scalar field $\phi$, we would expect to find terms in the Lagrangian
density of the form $\phi\Box^n\phi$. Any one such term would make a direct
contribution to the scalar self-energy function $\Pi^*(q^2)$ proportional to
$(q^2)^n$. If we were to include this contribution to all orders, but ignore
all other effects of non-renormalizable interactions, then the propagator
$\Delta'(q^2)=1/(q^2+m^2-\Pi^*(q^2))$ would not have the simple pole in $q^2$
at negative $q^2$ expected from the general arguments of Section 10.7, but
$n$ such poles (some of which may coincide), generally at complex values of
$q^2$. But if the non-renormalizable term $\phi\Box^n\phi$ has a coefficient
of order $M^{-2(n-1)}$, where $M\gg m$, then the extra poles are at $q^2$ of
order $M^2$, where it is illegitimate to ignore the infinite number of other
non-renormalizable interactions that must also appear in the Lagrangian. Thus
the appearance of higher-derivative terms in a general non-renormalizable
Lagrangian is not in conflict with the general principles underlying quantum
field theory that were used in Section 10.7. But by the same token, we also
cannot use higher-derivative terms to avoid ultraviolet divergences
altogether, as has been repeatedly proposed. A term
$M^{-2(n-1)}\phi\Box^n\phi$ in the Lagrangian density provides a cutoff at
momenta $q^2\approx M^2$, but at these momenta we cannot ignore all the other
non-renormalizable interactions that must be present.

Although highly suppressed, non-renormalizable interactions may be detectable
if they have effects that would otherwise be forbidden. For instance, we will
see in Section 12.5 that the symmetries of charge-conjugation and
space-inversion invariance are an automatic consequence of the structure of
the electromagnetic interactions that is imposed by gauge invariance, Lorentz
invariance, and renormalizability, but we can easily imagine
non-renormalizable terms that would violate these symmetries, such as an
electron electric dipole moment term
$\bar\psi\gamma_5[\gamma_\mu,\gamma_\nu]\psi F^{\mu\nu}$, or the Fermi
interaction $\bar\psi\gamma_5\gamma_\mu\psi\bar\psi\gamma^\mu\psi$. It is
widely believed today that the conservation of baryon and lepton number is
violated by very small effects of highly suppressed non-renormalizable
interactions. Another example of a detectable non-renormalizable interaction
is provided by gravitation. As mentioned before, gravitons have no
renormalizable interactions at all. But, of course, we detect gravitation,
because it has the special property that the gravitational fields of all the
particles in a macroscopic body add up coherently.

Although non-renormalizable theories involve an infinite number of free
parameters, they retain considerable predictive
power:~\hyperref[ref:12.12]{[12]} they allow us to calculate the non-analytic
parts of Feynman amplitudes, like the $\ln q$ and $q\ln q$ terms in the
one-dimensional examples at the beginning of the previous section. Such
calculations just reproduce the results required by the axiom of $S$-matrix
theory, that the $S$-matrix has only those singularities required by unitarity.

Paradoxically, it is just in the case where symmetry principles forbid
renormalizable interactions that non-renormalizable quantum field theories
prove the most useful. In such cases we can derive a useful perturbation theory
by expanding in powers of $k/M$. This has been worked out in detail for the
theory of low-energy pions~\hyperref[ref:12.12]{[12,}\hyperref[ref:12.13]{13]},
to be discussed in detail in Volume II, and the theory of low-energy
gravitons~\hyperref[ref:12.14]{[14]}. For a simpler example, consider the
theory of a real scalar field, satisfying the principle of invariance under
the field translation
\[
\phi(x)\longrightarrow\phi(x)+\epsilon
\]
with $\epsilon$ an arbitrary constant. This symmetry forbids any renormalizable
interactions or scalar mass, but it allows an infinite number of
non-renormalizable derivative interactions
\[
\mathcal L=-\frac12\partial_\mu\phi\partial^\mu\phi
-\frac{g}{4}(\partial_\mu\phi\partial^\mu\phi)^2-\cdots\,;
\]
where $g\approx M^{-4}$, and `$\cdots$' denotes terms with more derivatives or
fields. (For simplicity, it is assumed here that the theory also has a
symmetry under the reflection $\phi\to-\phi$.) According to the above
dimensional analysis, the graph for a general reaction in which all energies
and momenta are of order $k\ll M$ is suppressed by a factor $(k/M)^\nu$, where
\[
\nu=-\sum_i V_i\Delta_i=\sum_i V_i(d_i+n_i-4),
\]
with $n_i$ and $d_i$ the numbers of scalar fields and derivatives in an
interaction of type $i$, and $V_i$ the number of vertices for these
interactions in our graph. For $k\ll M$, the dominant contributions to any
process are those with the smallest value of $\nu$. The formula for $\nu$ can
be put in a more useful form by using the familiar topological identities for
a connected graph:
\[
\sum_i V_i=I-L+1,\qquad \sum_i V_i n_i=2I+E,
\]
where $I$, $E$, and $L$ are the numbers of internal lines, external lines,
and loops in our graph. Combining these relations gives
\[
\nu=2E-4+4L+\sum_i V_i(d_i-n_i),
\]
Now, the field translation symmetry requires that every field must be
accompanied with at least one derivative, so the quantity $d_i-n_i$ as well as
$L$ is non-negative for all interactions. Thus for a given process (that is,
a fixed number $E$ of external lines) the dominant terms will be those
constructed solely from \emph{tree} graphs (i.e., $L=0$), and interactions
with the minimum number $d_i=n_i$ of derivatives. That is, in leading order
we can take the Lagrangian density to depend only on \emph{first} derivatives
of the field. Higher-order corrections may involve loops and/or interactions
with more derivatives on some fields. But to any given order $\nu$ in $k/M$,
we need only consider a finite number of graphs, those with
$L\leq(4-2E+\nu)/4$, and only a finite number of interaction types.

For instance, scalar--scalar scattering is given in leading order by the
one-vertex tree graph calculated using the interaction
$-g(\partial_\mu\phi\partial^\mu\phi)^2$ in first order. According to our
formula for $\nu$, the leading correction, suppressed at low energy by a
factor $(k/M)^2$, arises from another single-vertex tree graph, produced by an
interaction with two additional derivatives of the
form\footnote{In accordance with the remarks of Section 7.7, we are excluding
interactions involving $\Box\phi$, because the field equation for $\phi$ can
be used to express such interactions in terms of the others.}
$\partial_\mu\partial_\nu\phi\partial^\mu\partial^\nu\phi
\partial_\rho\phi\partial^\rho\phi$. The next corrections, suppressed at low
energy by two further factors of $k/M$, arise both from the one-loop diagram
of Figure~\ref{fig:weinberg-12-4} (including permutations of external lines),
calculated using only the interaction
$-g(\partial_\mu\phi\partial^\mu\phi)^2$, and also from tree graphs with a
single vertex arising from a quartic interaction with eight derivatives,
whose couplings contain infinite parts that cancel the ultraviolet divergence
from the loop graph.\footnote{These are the only ultraviolet divergences
encountered in one-loop graphs if we use dimensional regularization. For
other methods of regularization there are also quartic and quadratic
divergences, which are cancelled by counterterms in the four-scalar
interactions with four or six derivatives.} The loop graph also yields finite
terms in the scattering amplitude proportional to terms like
$s^4\ln s+t^4\ln t+u^4\ln u$, $s^2t^2\ln u+t^2u^2\ln s+u^2s^2\ln t$, etc.,
with calculable coefficients proportional to $g^2$. These finite terms simply
represent the correction to the lowest-order scattering amplitude needed to
ensure the unitarity of the $S$-matrix, but perturbative quantum field theory
is by far the easiest way of calculating them.

% Figure 12.4, physical PDF p. 546 (printed p. 523).
\begingroup
\renewcommand{\thefigure}{12.\arabic{figure}}
\setcounter{figure}{3}
\begin{figure}[H]
\centering
\begin{tikzpicture}[scalar/.style={draw,line width=0.75pt,dashed}]
  \coordinate (L) at (-1.65,0);
  \coordinate (R) at (1.65,0);
  \draw[scalar] (-3.25,1.45) -- (L);
  \draw[scalar] (-3.25,-1.45) -- (L);
  \draw[scalar] (R) -- (3.25,1.45);
  \draw[scalar] (R) -- (3.25,-1.45);
  \draw[scalar] (L) .. controls (-0.70,1.00) and (0.70,1.00) .. (R);
  \draw[scalar] (L) .. controls (-0.70,-1.00) and (0.70,-1.00) .. (R);
  \fill (L) circle (1.15pt);
  \fill (R) circle (1.15pt);
\end{tikzpicture}
\caption{One-loop diagram for scalar--scalar scattering in the theory with
derivative quadrilinear interactions.}
\label{fig:weinberg-12-4}
\end{figure}
\endgroup


Although non-renormalizable theories can provide useful expansions in powers
of energy, they inevitably lose all predictive power at energies of the order
of the common mass scale $M$ that characterizes the various couplings. If we
were to take these expansions literally, the results for $S$-matrix elements
would violate unitarity bounds for $E\gg M$. There seem to be just two
possibilities about what happens at such energies. One is that the growing
strength of the effects of the non-renormalizable couplings somehow saturates,
avoiding any conflict with unitarity~\hyperref[ref:12.15]{[15]}. The other is
that new physics of some sort enters at the scale $M$. In this case, the
non-renormalizable theories that describe nature at energies $E\ll M$ are just
\emph{effective field theories} rather than truly fundamental theories.

Probably the earliest example of an effective field theory was derived in the
1930s by Euler \emph{et al.}~\hyperref[ref:12.16]{[16]}, as a theory of
low-energy photon--photon interactions. (See Section 1.3.) In effect, they
calculated the contribution to photon--photon scattering of Feynman diagrams
such as Figure~\ref{fig:weinberg-12-5}, and found that at energies much less
than $m_e$ the scattering of light by light was the same as would be
calculated with an effective Lagrangian
\[
\begin{split}
\mathcal L_{\mathrm{eff}}=\frac{2\alpha^2}{45m_e^4}
\left[(\mathbf E^2-\mathbf B^2)^2+7(\mathbf E\cdot\mathbf B)^2\right]
\qquad\qquad\\
{}+\text{higher orders in }\frac{eE}{m_e^2}\ \&\ \frac{eB}{m_e^2}
\end{split}
\]

% Figure 12.5, physical PDF p. 547 (printed p. 524).
\begingroup
\renewcommand{\thefigure}{12.\arabic{figure}}
\setcounter{figure}{4}
\begin{figure}[H]
\centering
\begin{tikzpicture}[
  fermion/.style={line width=0.75pt,postaction={decorate},
    decoration={markings,mark=at position 0.55 with
      {\arrow{Stealth[length=5pt,width=4pt]}}}},
  photon/.style={line width=0.7pt,decorate,
    decoration={snake,amplitude=1.25pt,segment length=6pt}}
]
  \coordinate (TL) at (-1.35,1.35);
  \coordinate (TR) at (1.35,1.35);
  \coordinate (BR) at (1.35,-1.35);
  \coordinate (BL) at (-1.35,-1.35);
  \draw[fermion] (TL) -- (TR);
  \draw[fermion] (TR) -- (BR);
  \draw[fermion] (BR) -- (BL);
  \draw[fermion] (BL) -- (TL);
  \draw[photon] (-2.15,2.75) -- (TL);
  \draw[photon] (2.15,2.75) -- (TR);
  \draw[photon] (BR) -- (2.15,-2.75);
  \draw[photon] (BL) -- (-2.15,-2.75);
\end{tikzpicture}
\caption{Diagram for photon--photon scattering, whose effect at low energy can
be calculated from the effective Lagrangian of Euler \emph{et
al.}~\hyperref[ref:12.16]{[16]}. Straight lines are electrons; wavy lines are
photons.}
\label{fig:weinberg-12-5}
\end{figure}
\endgroup


Euler \emph{et al.} used this effective Lagrangian only in the tree
approximation, to calculate the leading terms in photon interaction matrix
elements. It was not until much later that such Lagrangians, though
non-renormalizable, were used beyond the tree
approximation~\hyperref[ref:12.12]{[12,}\hyperref[ref:12.17]{17]}.

In modern jargon, we say that in deriving this Lagrangian the electron is
`integrated out', because in the one-loop approximation we have
\[
\exp\left(i\int\mathcal L_{\mathrm{eff}}(\mathbf E,\mathbf B)\,d^4x\right)
=\int\left[\prod_x d\psi_e(x)\right]
\exp\left(i\int\mathcal L_{\mathrm{QED}}(\psi_e,A)\,d^4x\right).
\]
A more general procedure is simply to write down the most general
non-renormalizable effective Lagrangian, use it to calculate various
amplitudes as an expansion in energies and momenta, and then choose the
constants in the effective Lagrangian by matching the results it gives for
these amplitudes to those derived from the underlying theory.

We will encounter effective field theories again, especially in considering
broken symmetries in Volume II. As we shall see, effective field theories are
useful even where they cannot be derived from an underlying theory, either
because the theory is unknown, or because its interactions are too strong to
allow the use of perturbation theory. Indeed, even if we knew nothing about
the properties of charged particles, the scattering of photons at
sufficiently low energy would have to be described by an effective Lagrangian
consisting of the terms $(\mathbf E^2-\mathbf B^2)^2$ and
$(\mathbf E\cdot\mathbf B)^2$, because these are the unique quartic Lorentz-
and gauge-invariant terms with no derivatives acting on $\mathbf E$ and
$\mathbf B$. Terms with such derivatives would be suppressed at low photon
energies $E$ by additional factors of $E/M$, where $M$ is some typical mass of
the charged particles that are being integrated out. We can go further: we
shall see that effective field theories are useful even where the light
particles they describe are not present in the underlying theory at all, but
composites of the heavy particles that are integrated out. The underlying
theory might not even be a field theory at all---the problem of incorporating
gravitation has led many theorists to believe that, in fact, it is a string
theory. But wherever an effective field theory comes from, it is inevitably a
non-renormalizable theory.
